\section{Experimental Results}
This section is intentionally written as a \emph{draft experimental narrative}. The setup description is taken from the internal HW-SW partitioning draft and the repository code, while the numeric evidence comes from result files already present in the workspace. The section therefore demonstrates how the final paper could be structured, but the numbers should be rerun before any real submission.

\subsection{Experimental Setting}
The internal draft describes task graphs derived from real neural-network models at tensor-operation granularity. The graph topology is taken from those models, while node and edge costs are assigned synthetically. In particular:
\begin{itemize}
\item software costs are sampled from a uniform range;
\item hardware costs are sampled around a scaled software cost using a factor $k$ and variance parameter $l$;
\item hardware area is sampled independently up to a chosen maximum;
\item communication costs are sampled relative to the largest software execution time using a scale factor $\mu$.
\end{itemize}
This setup is also consistent with \texttt{TaskGraph.load\_graph\_from\_pydot}, which generates software time, hardware time, area, and communication values from a loaded graph topology.

For the differentiable GNN methods, the most relevant evaluation metrics are:
\begin{itemize}
\item \textbf{Queue makespan}, produced by the repository's queue-style simulator;
\item \textbf{Legacy LP makespan}, produced by the older scheduling backend;
\item \textbf{Partition cost}, which is useful for diagnosis but not sufficient by itself.
\end{itemize}
The main comparison in this draft is between \texttt{diff\_gnn} and \texttt{diff\_gnn\_order}. Classical methods such as GL25, GCPS, and PSO remain valuable reference points, but the purpose here is to understand what the ordering extension contributes beyond the placement-only differentiable model.

\subsection{Repository-Derived Example on the Eleven-Node Figure-3 Graph}
The repository contains a small but fully analyzable case study in the file \texttt{fig3\_area18\_paper\_dls\_eval\_summary.csv}. The graph has eleven tasks and an area budget of approximately $18.0$ units. Table~\ref{tab:fig3_eval} shows the recorded results for MIP, the baseline differentiable model, and the ordered variant.

\begin{table}[t]
\centering
\caption{Repository-derived results on the eleven-node demonstration graph.}
\label{tab:fig3_eval}
\small
\begin{tabular}{lrrrrr}
\toprule
Method & Mksp. & Part. & Area & Cross & Comm. \\
\midrule
MIP & 34 & 66 & 18.0 & 6 & 12 \\
Diff-GNN+Order & 35 & 69 & 15.0 & 6 & 14 \\
Diff-GNN & 37 & 66 & 16.0 & 5 & 8 \\
\bottomrule
\end{tabular}
\end{table}

This table gives a useful, nontrivial story. The ordered variant improves queue makespan from $37$ to $35$ relative to the baseline, moving closer to the MIP solution at $34$. However, the improvement does not come from a lower partition cost; in fact the recorded partition cost is slightly higher. This is exactly the kind of result one would expect in a schedule-sensitive partitioning problem. A partition can look worse under a static cost proxy yet produce a better actual schedule because it better matches precedence and resource interactions.

Another useful observation is that the ordered model does not simply place more nodes in hardware. In this example it uses \emph{less} hardware area than the baseline while improving makespan. That suggests the gain comes from \emph{where} hardware is used and how the resulting partition interacts with the software schedule, not from saturating the available area budget.

\subsection{Mixed Behavior Across Sweep Configurations}
A stronger paper should also show cases where ordering does \emph{not} help. The repository already contains such evidence in \texttt{fig3\_sweep\_ranked\_2026-02-19.csv}. Two cases are especially informative:
\begin{itemize}
\item Under \texttt{config\_fig3\_sweep\_v04\_minimal}, the baseline records queue makespan $39$ while the ordered variant records $35$, indicating a clear benefit from the extra order model.
\item Under \texttt{config\_fig3\_sweep\_v02\_noreg}, the baseline records queue makespan $35$ while the ordered variant records $37$, indicating a regression.
\end{itemize}

The same sweep also shows the runtime cost of the ordered model. For example, under \texttt{v04\_minimal}, the baseline requires about $13.8$s whereas the ordered model requires about $26.8$s. Under \texttt{v08\_order\_strong}, the ordered model takes more than $64$s and still does not beat the best baseline setting. The lesson is that ordering changes both the optimization landscape and the wall-clock budget. A careful final study should therefore report quality and runtime together.

These mixed results are not surprising after inspecting the code. The ordered model introduces extra temperatures, soft permutations, and pairwise-order computations. At the same time, the current default configuration keeps \texttt{use\_hw\_ordering=False}, so only the software-order head is fully active in the loss. The method is therefore best viewed as a partially activated order-aware extension whose benefit depends on graph structure and regularization.

\subsection{Aggregate Comparison on a Fifteen-Instance Random Benchmark}
The repository also contains an aggregate summary in \texttt{shared\_random\_15\_20260226\_070718\_summary.csv}. Table~\ref{tab:random15} reports the values for the two GNN variants and two older baselines.

\begin{table}[t]
\centering
\caption{Aggregate summary over fifteen random instances from repository outputs.}
\label{tab:random15}
\small
\begin{tabular}{lrrr}
\toprule
Method & Mean Mksp. & Median Mksp. & Best Cnt. \\
\midrule
Diff-GNN+Order & 4624.54 & 5210.00 & 13 \\
Diff-GNN & 4921.82 & 5423.63 & 4 \\
GL25 & 6653.44 & 5978.47 & 0 \\
GCPS & 6908.76 & 8369.32 & 0 \\
\bottomrule
\end{tabular}
\end{table}

Here the ordered model is more clearly favorable. It improves mean makespan by roughly $6\%$ and median makespan by about $4\%$ relative to the placement-only differentiable model, and it wins thirteen out of fifteen recorded cases. These numbers are not enough for a polished claim because the exact benchmark generation and repetition protocol should still be rerun from scratch. Still, they illustrate how a final paper can argue that the ordering extension is valuable on average even if it is not uniformly dominant on every single small instance.

This aggregate table also suggests a broader narrative. Both differentiable GNN variants outperform older heuristic baselines in the summary, which supports the general premise that graph-based learned surrogates are promising for large combinatorial partitioning spaces. Within that family, the order-aware variant appears to offer a second layer of improvement by modeling schedule structure rather than only placement.

\subsection{Interpretation}
Taken together, the repository results suggest three defensible conclusions for a final manuscript.

First, schedule-aware partitioning should be evaluated by actual schedule metrics, not by partition cost alone. The eleven-node example already shows a partition-cost/makespan mismatch.

Second, the ordered extension has a meaningful upside. On the aggregate fifteen-instance summary it improves both mean and median makespan, and on at least some small controlled settings it closes part of the gap between the baseline differentiable model and MIP.

Third, ordering is not free. It increases runtime, enlarges the hyperparameter space, and can regress under weak regularization or mismatched settings. This makes an ablation study essential. In a final experimental section, the following tables would be particularly useful:
\begin{itemize}
\item with versus without ordering under identical feature and postprocess settings;
\item with software ordering only versus software+hardware ordering;
\item queue metric versus legacy LP metric;
\item no postprocess versus DLS versus DLS+LSSP.
\end{itemize}

\subsection{Recommended Final Experimental Structure}
For later manual rewriting, a full version of this section could be organized as follows:
\begin{enumerate}
\item \textbf{Setup.} Describe graph sources, the $(k,l,\mu,\alpha)$ cost model, seeds, and train/eval metrics.
\item \textbf{Method ablations.} Compare \texttt{diff\_gnn} against \texttt{diff\_gnn\_order}, then switch on or off paper features, learned edge weights, DLS, and hardware ordering.
\item \textbf{Comparisons to baselines.} Add GCPS, GL25, PSO-family methods, and MIP where tractable.
\item \textbf{Runtime analysis.} Report wall-clock time and memory use.
\item \textbf{Schedule diagnostics.} Include cross-edge counts, communication delay, area usage, and critical-path behavior for representative graphs.
\end{enumerate}

That structure is consistent with the repository and with the two reference PDFs. It also gives the user a practical template for rewriting the experiments later in their own wording and with their own rerun data.
